\section{Introduction}
\label{sec:intro}
The Advanced Encryption Standard (AES) is a central component of systems that protect stored and transmitted data~\cite{fips197}. Its computational cost matters in applications ranging from bulk storage encryption to services processing short messages. In these settings, encryption shares processor time and memory bandwidth with other application functions. Understanding its performance therefore requires attention both to the cipher implementation and to the software through which an application invokes it.

Modern processors provide dedicated cryptographic instructions, while libraries select among implementations according to processor capabilities and software configuration. However, applications also initialize cipher contexts, allocate output buffers, convert data representations and transfer control between programming languages. A measurement taken within a native encryption routine excludes some of this work. A measurement around a higher-level call includes more of it. Both are useful, but each describes a different part of the computation.

This distinction matters when evaluating native libraries through Python interfaces. Libraries implementing the same AES mode may perform different preparation and data-handling operations. Their measured difference can consequently reflect cipher implementation, interface design and timer placement. These contributions must be distinguished before using a benchmark to select a library or estimate application capacity.

This paper addresses the following research question: \emph{how do the selected AES implementation and the scope of the timed operation affect measured CPU throughput, and which conclusions remain useful for application design?} We investigate this question through a retrospective analysis of validated measurements from a Windows x86-64 laptop and an Oracle Cloud Arm64 instance. AES-128 in counter (CTR) mode provides the main workload. Separate Windows experiments with cipher block chaining (CBC) and Galois/Counter Mode (GCM) extend the analysis to key length and authenticated processing~\cite{sp80038a,sp80038d}.

The paper makes three contributions. First, it quantifies the differences between the evaluated OpenSSL and PyCryptodome implementations on each host, including an OpenSSL capability-mask experiment on Arm64. Second, it shows how the relative performance observed on Windows changes when timing expands from encryption processing to the complete instrumented benchmark call. Third, it relates these observations to the effects of key length and message size in separate CBC and GCM workloads. Together, these contributions identify the information needed to interpret an encryption throughput figure and assess its relevance to a software service.

The analysis uses existing experimental records; no additional measurements were collected for this paper. Comparisons are made within each host and workload, with the timing differences stated explicitly. Section~\ref{sec:related} places the study in the context of AES implementation and benchmarking research. Section~\ref{sec:method} describes the experimental setup and estimators. Section~\ref{sec:results} presents the results and their interpretation, followed by broader implications, limitations and future work in Section~\ref{sec:discussion}. Section~\ref{sec:conclusion} concludes the paper.
